\chapter{Urn Model}
\label{ch:3}
\section{Markov chains}
\label{sec:3-1}
In the discrete-time Ehrenfest chain, the state $j$ is the number of balls in one of two urns. Moving one uniformly chosen ball gives
\[
(Pf)(j)=\frac{N-j}{N}f(j+1)+\frac jN f(j-1).
\]
The Krawtchouk difference equation gives
\[
P K_n(\cdot;N)=\left(1-\frac{2n}{N}\right)K_n(\cdot;N).
\]
Here $N\geq1$, and terms outside the state space have zero transition coefficient.

\section{Initial and invariant distributions}
\label{sec:3-2}
The invariant distribution is $w_N(j)=2^{-N}\binom Nj$. Detailed balance follows from
\[
w_N(j)\frac{N-j}{N}=w_N(j+1)\frac{j+1}{N}.
\]
The non-lazy discrete-time chain has period two; invariance should not be confused with convergence to equilibrium from every initial distribution.


\section*{Connection to the supplied article sources}
Section~4 (``Ehrenfest urn model'') of the supplied
\texttt{KravFeinsKocik.tex} belongs to this chapter.
The urn state is the Hamming weight of a binary configuration: flipping
a uniformly selected coordinate moves from $j$ to $j+1$ with probability
$(N-j)/N$ and to $j-1$ with probability $j/N$, for $N\geq1$.
In the book's column-vector convention the transition matrix is
$A_N/N$. The tensor-space realization is discussed in Chapter~\ref{ch:9}.
